\section{条件生成与 CFG}

\subsection{从无条件生成到条件生成}

真正有用的生成模型几乎都不是无条件的。我们通常想要的是“给定类别、文本、草图、边界框、掩码以后再生成”。在 rectified flow 或 diffusion 里，这很自然：把条件 $y$ 作为额外输入即可。

\begin{figure}[H]
\centering
\includegraphics[width=0.95\textwidth]{slides-images/slide-061.jpg}
\caption{条件 rectified flow：模型不仅要知道噪声状态，还要知道条件 $y$\protect\footnotemark}
\end{figure}
\footnotetext{Slides 第62页。}

\begin{knowledgebox}{条件生成的真实需求}
工程上，条件不是“锦上添花”，而是生成系统是否能落地的关键：
\begin{itemize}
    \item 文本生成图像需要 prompt
    \item 图像编辑需要掩码或参考图
    \item 视频生成需要时间一致性和语义约束
    \item 产品级系统必须允许用户控制结果
\end{itemize}
\end{knowledgebox}

\subsection{Classifier-free guidance 的核心思想}

CFG 的做法非常巧：训练时随机把条件 $y$ 丢掉，让同一个模型同时学会 conditional 和 unconditional 两种模式。于是推理时我们可以分别得到

\[
v_{\emptyset}=f_\theta(x_t,\varnothing,t), \qquad v_y=f_\theta(x_t,y,t).
\]

然后用线性组合控制“听话程度”：

\[
v_{\text{cfg}}=(1+w)v_y-wv_{\emptyset}.
\]

\begin{figure}[H]
\centering
\includegraphics[width=0.95\textwidth]{slides-images/slide-065.jpg}
\caption{CFG 的两条向量场：一个是无条件方向，一个是带条件方向\protect\footnotemark}
\end{figure}
\footnotetext{Slides 第66页。}

\begin{figure}[H]
\centering
\includegraphics[width=0.95\textwidth]{slides-images/slide-066.jpg}
\caption{CFG 的线性组合：通过放大条件向量场，把采样方向更明显地拉向条件分布\protect\footnotemark}
\end{figure}
\footnotetext{Slides 第67页。}

\begin{importantbox}{CFG 的直觉}
可以把 $v_{\emptyset}$ 理解成“往一般数据流形走”，把 $v_y$ 理解成“往满足条件的子流形走”。CFG 不是重新训练一个模型，而是通过线性外推，把采样方向更强地推向条件分布。
\end{importantbox}

\subsection{为什么 CFG 如此重要}

CFG 之所以几乎成了现代扩散模型标配，是因为它同时解决了两个问题：
\begin{enumerate}
    \item 不需要额外训练一个独立 classifier。
    \item 采样时可以直接控制 condition strength。
\end{enumerate}

\begin{figure}[H]
\centering
\includegraphics[width=0.95\textwidth]{slides-images/slide-070.jpg}
\caption{CFG 的来源：早期方法需要额外分类器来估计 $\nabla_x \log p(y\mid x)$，CFG 则把这件事并入同一个生成模型\protect\footnotemark}
\end{figure}
\footnotetext{Slides 第71页。}

\begin{warningbox}{CFG 的代价}
CFG 的代价也很直接：采样时通常要跑两次模型，一次 conditional，一次 unconditional，所以推理成本会翻倍。更大的 guidance scale 往往更“听话”，但也更容易牺牲多样性。
\end{warningbox}

\begin{figure}[H]
\centering
\includegraphics[width=0.95\textwidth]{slides-images/slide-071.jpg}
\caption{CFG 的实战结论：它非常重要，也确实会增加采样成本\protect\footnotemark}
\end{figure}
\footnotetext{Slides 第72页。}

